\documentclass[11pt,a4paper]{article}

\usepackage[utf8]{inputenc}
\usepackage[T1]{fontenc}
\usepackage[french]{babel}
\usepackage{amsmath,amssymb,amsfonts}
\usepackage{geometry}
\usepackage{graphicx}
\usepackage{booktabs}
\usepackage{hyperref}
\usepackage{array}
\usepackage{xcolor}
\usepackage{float}

\geometry{margin=2.5cm}
\graphicspath{{figures/}}

\title{
\textbf{Tenseur énergie--impulsion émergent depuis l'intrication}\\
\large Source modulaire et réponse de courbure dans Bottom-Up Quantum Gravity
}

\author{Farid Hamdad}
\date{\today}

\begin{document}

\maketitle

\begin{abstract}
Paper~18 étudie la troisième flèche nécessaire à la limite d'Einstein
effective de Bottom-Up Quantum Gravity :
\[
\delta W_{\rm loc}
\longrightarrow
T_{\mu\nu}^{\rm ent}.
\]
Paper~16 a renforcé la convergence spectrale
\(L_N\to\Delta_g\), et Paper~17 a identifié un proxy moyen de Ricci à partir
de la courbure d'Ollivier--Ricci. Le présent papier traite le côté source :
comment une perturbation locale du graphe d'intrication peut-elle jouer le
rôle d'une source énergie--impulsion effective ?

Nous construisons un analogue graphe de la première loi modulaire. Sur une
région \(A\), nous définissons
\[
\rho_A=
\frac{(L_A+\mu I)^{-1}}{\mathrm{Tr}(L_A+\mu I)^{-1}},
\qquad
S_A=-\mathrm{Tr}(\rho_A\log\rho_A),
\qquad
K_A=-\log\rho_A.
\]
Après une perturbation locale de \(W_{ij}\), nous vérifions
\[
\delta S_A^{\rm graph}
\simeq
\delta\langle K_A^{\rm graph}\rangle.
\]
Sur le tore plat et la sphère, la pente du fit vaut respectivement
\(0.989481\) et \(0.989718\), avec \(R^2\simeq0.9966\).

Nous montrons ensuite que la variation modulaire prédit la réponse de
courbure locale :
\[
|\delta\langle K_A\rangle|
\longrightarrow
\langle|\Delta\kappa|\rangle_{\rm near}.
\]
Le coefficient de détermination vaut \(R^2=0.967229\) sur le tore plat et
\(R^2=0.982252\) sur la sphère. Ces résultats valident numériquement la chaîne
\[
\delta W_{\rm loc}
\to
\delta S_A
\simeq
\delta\langle K_A\rangle
\to
\delta\kappa(r),
\]
et fournissent une première source modulaire effective pour la courbure.
\end{abstract}

\tableofcontents

\section{Introduction}

Dans Bottom-Up Quantum Gravity, la géométrie et la gravité émergent de la
structure d'intrication d'un état quantique global. L'objet fondamental est
la matrice d'information mutuelle
\[
W_{ij}=I(i:j).
\]

Les papiers précédents ont établi deux briques géométriques. Paper~16 a
montré que le Laplacien d'intrication rescalé reconstruit le bas spectre du
Laplacien de Laplace--Beltrami :
\[
c_N L_N\to-\Delta_g.
\]
Paper~17 a montré que la courbure d'Ollivier--Ricci porte un signal de Ricci
moyen :
\[
\frac{\kappa^{OR}}{\epsilon}
\simeq
B_N+C_N R_{\mu\nu}u^\mu u^\nu.
\]

Il reste à comprendre le côté source de l'équation d'Einstein effective :
\[
T_{\mu\nu}^{\rm ent}.
\]

Paper~18 pose donc la question :
\[
\boxed{
\text{Une perturbation locale de }W_{ij}
\text{ définit-elle une source effective de courbure ?}
}
\]

La chaîne visée est :
\[
\delta W_{\rm loc}
\to
\delta S_A
\simeq
\delta\langle K_A\rangle
\to
\delta\kappa(r)
\to
T_{\mu\nu}^{\rm ent}.
\]

\section{Héritage de Paper 8 et Paper 15}

Paper~8 avait déjà proposé une construction effective :
\[
T_{\mu\nu}^{\rm eff}
=
T_{\mu\nu}^{\rm matter}[\delta W_{\rm loc}]
+
T_{\mu\nu}^{\rm ent}[d_s].
\]
Le meilleur proxy source présentait une corrélation de Spearman
\[
\rho=0.741,
\qquad
p=1.84\times10^{-4}
\]
avec les fluctuations de courbure.

Paper~15 Step~E a ensuite montré qu'une perturbation radiale de l'intrication
produit une réponse de courbure localisée :
\[
\delta W_{\rm loc}
\to
\delta\kappa(r).
\]
Sur la sphère, le test donnait
\[
{\rm Spearman}(\phi_{\rm edge},|\Delta\kappa|)
=
0.752,
\]
et un ratio near/far
\[
11.51.
\]

Paper~18 poursuit cette direction, mais introduit explicitement la variation
modulaire comme source intermédiaire.

\section{Première loi modulaire}

La première loi modulaire affirme :
\[
\delta S_A
=
\delta\langle K_A\rangle,
\]
où
\[
K_A=-\log\rho_A.
\]

Dans une théorie quantique des champs relativiste, le Hamiltonien modulaire
local peut être relié au tenseur énergie--impulsion :
\[
\delta\langle K_A\rangle
\sim
\int_A \xi^\mu\delta T_{\mu\nu}\,d\Sigma^\nu.
\]

Dans BuP, l'objectif est de construire un analogue discret :
\[
\delta S_A[W]
\simeq
\delta\langle K_A[W]\rangle,
\]
puis d'interpréter
\[
\delta\langle K_A\rangle
\]
comme une source effective de courbure.

\section{Construction graphe de \(\rho_A\), \(S_A\) et \(K_A\)}

Sur une région \(A\) du graphe, on restreint le Laplacien :
\[
L_A=L|_A.
\]

On définit une matrice densité effective :
\[
\rho_A
=
\frac{(L_A+\mu I)^{-1}}
{\mathrm{Tr}(L_A+\mu I)^{-1}}.
\]

Le paramètre \(\mu\) joue le rôle d'une régularisation infrarouge. Il garantit
que l'inverse est bien défini.

L'entropie de graphe est :
\[
S_A
=
-\mathrm{Tr}(\rho_A\log\rho_A).
\]

Le Hamiltonien modulaire discret est :
\[
K_A=-\log\rho_A.
\]

Après une perturbation locale du graphe,
\[
W_{ij}
\to
W'_{ij},
\]
on obtient
\[
\rho_A\to\rho'_A.
\]

On mesure alors :
\[
\delta S_A
=
S_A(W')-S_A(W),
\]
et :
\[
\delta\langle K_A\rangle
=
\mathrm{Tr}\left[(\rho'_A-\rho_A)K_A\right].
\]

La première loi modulaire prédit :
\[
\delta S_A
\simeq
\delta\langle K_A\rangle.
\]

\section{Perturbation locale d'intrication}

La perturbation utilisée est radiale :
\[
\phi_i
=
\exp\left[
-\frac{d(i,\mathrm{source})^2}{2\sigma^2}
\right].
\]

Les poids sont modifiés selon :
\[
W'_{ij}
=
W_{ij}
\left[
1+s\frac{\phi_i+\phi_j}{2}
\right].
\]

Le paramètre \(s\) contrôle l'intensité de la source. Les tests utilisent :
\[
s\in\{-0.20,-0.10,-0.05,0.05,0.10,0.20\}.
\]

\section{Résultat v1 : loi modulaire sur graphe}

Le premier test mesure la relation entre
\[
\delta S_A
\]
et
\[
\delta\langle K_A\rangle.
\]

Sur le tore plat, le fit linéaire donne :
\[
\delta S_A
=
-0.000068
+
0.989481\,\delta\langle K_A\rangle.
\]
On obtient :
\[
R^2=0.996590,
\qquad
r_{\rm Pearson}=0.998293.
\]

Sur la sphère :
\[
\delta S_A
=
-0.000081
+
0.989718\,\delta\langle K_A\rangle.
\]
On obtient :
\[
R^2=0.996592,
\qquad
r_{\rm Pearson}=0.998295.
\]

L'erreur relative moyenne vaut environ :
\[
0.0533
\]
sur les deux géométries.

Ainsi :
\[
\boxed{
\delta S_A^{\rm graph}
\simeq
\delta\langle K_A^{\rm graph}\rangle.
}
\]

\section{Résultat v2 : source modulaire et réponse de courbure}

La version v2 mesure simultanément la réponse modulaire et la réponse de
courbure :
\[
\Delta\kappa_{ij}
=
\kappa_{ij}^{\rm after}
-
\kappa_{ij}^{\rm before}.
\]

On compare trois zones :

\[
{\rm near},
\qquad
{\rm mid},
\qquad
{\rm far}.
\]

La réponse de courbure est localisée. Le ratio
\[
\frac{
\langle|\Delta\kappa|\rangle_{\rm near}
}{
\langle|\Delta\kappa|\rangle_{\rm far}
}
\]
vaut :
\[
6.293
\]
sur le tore plat, et :
\[
9.534
\]
sur la sphère.

Le résultat principal est que l'amplitude modulaire prédit l'amplitude de la
réponse de courbure proche :
\[
|\delta\langle K_A\rangle|
\longrightarrow
\langle|\Delta\kappa|\rangle_{\rm near}.
\]

Sur le tore plat :
\[
R^2=0.967229,
\qquad
r_{\rm Pearson}=0.983478.
\]

Sur la sphère :
\[
R^2=0.982252,
\qquad
r_{\rm Pearson}=0.991086.
\]

La relation signée est également quasi parfaite, à une convention de signe
près :
\[
R^2=0.998336,
\qquad
r_{\rm Pearson}=-0.999168
\]
sur le tore plat, et :
\[
R^2=0.998857,
\qquad
r_{\rm Pearson}=-0.999428
\]
sur la sphère.

Ces résultats montrent que :
\[
\boxed{
\delta\langle K_A\rangle
\text{ agit comme une source effective de la réponse de courbure.}
}
\]

\section{Résumé des résultats numériques}

\[
\begin{array}{lllll}
\toprule
{\rm Test} & {\rm Géométrie} & {\rm Quantité} & {\rm Valeur} & {\rm Statut}\\
\midrule
{\rm v1} & T^2 & {\rm pente}\ \delta S/\delta K & 0.989481 & positif\\
{\rm v1} & T^2 & R^2(\delta S,\delta K) & 0.996590 & positif\\
{\rm v1} & S^2 & {\rm pente}\ \delta S/\delta K & 0.989718 & positif\\
{\rm v1} & S^2 & R^2(\delta S,\delta K) & 0.996592 & positif\\
{\rm v2} & T^2 & {\rm near/far} & 6.293007 & positif\\
{\rm v2} & T^2 & R^2(|\delta K|,|\Delta\kappa|_{\rm near}) & 0.967229 & positif\\
{\rm v2} & S^2 & {\rm near/far} & 9.533529 & positif\\
{\rm v2} & S^2 & R^2(|\delta K|,|\Delta\kappa|_{\rm near}) & 0.982252 & positif\\
\bottomrule
\end{array}
\]

\section{Interprétation}

Les résultats de Paper~18 valident la chaîne :
\[
\delta W_{\rm loc}
\longrightarrow
\delta S_A
\simeq
\delta\langle K_A\rangle
\longrightarrow
\delta\kappa(r).
\]

La première relation montre que la variation d'entropie de graphe est bien
capturée par la variation modulaire. La seconde montre que cette variation
modulaire contrôle la réponse de courbure.

Le résultat n'est pas encore un tenseur énergie--impulsion complet. Il établit
un proxy scalaire de source :
\[
\delta\langle K_A\rangle,
\]
qui se comporte comme une source effective de courbure.

Dans la limite continue, l'objectif devient :
\[
\delta\langle K_A\rangle
\to
\int_A \xi^\mu \delta T_{\mu\nu}^{\rm ent}d\Sigma^\nu.
\]

\section{Lien avec l'équation d'Einstein effective}

Les résultats de Papers~16--18 fournissent les trois briques nécessaires :

\[
L_N\to\Delta_g,
\]

\[
\kappa^{OR}
\to
R_{\mu\nu}u^\mu u^\nu,
\]

\[
\delta W_{\rm loc}
\to
\delta\langle K_A\rangle
\to
\delta\kappa(r).
\]

Paper~19 pourra donc assembler :
\[
\frac{\delta S_{\rm BuP}}{\delta W_{ij}}=0
\quad
\Longrightarrow
\quad
G_{\mu\nu}
+
\Lambda_{\rm ent}g_{\mu\nu}
=
8\pi G_{\rm eff}T_{\mu\nu}^{\rm ent}
+
\mathcal{H}_{\mu\nu}.
\]

Paper~18 fournit la brique source :
\[
T_{\mu\nu}^{\rm ent}.
\]

\section{Limites}

Plusieurs limites doivent être précisées.

Premièrement, le test construit un proxy de matrice densité depuis le
Laplacien de graphe :
\[
\rho_A=
\frac{(L_A+\mu I)^{-1}}{\mathrm{Tr}(L_A+\mu I)^{-1}}.
\]
Ce n'est pas encore une vraie matrice densité quantique réduite \(\rho_A\).

Deuxièmement, le résultat est scalaire et modulaire. Il ne reconstruit pas
encore toutes les composantes de \(T_{\mu\nu}\).

Troisièmement, les géométries testées sont contrôlées et synthétiques. Il
reste à tester des graphes d'information mutuelle issus d'états quantiques.

Quatrièmement, la relation signée dépend de la convention de perturbation.

Cinquièmement, la dérivation continue via le théorème de
Bisognano--Wichmann reste à formaliser.

\section{Conclusion}

Paper~18 établit que la variation modulaire du graphe d'intrication agit comme
une source effective de courbure.

Le premier test valide :
\[
\delta S_A^{\rm graph}
\simeq
\delta\langle K_A^{\rm graph}\rangle.
\]

Le second test montre :
\[
|\delta\langle K_A\rangle|
\to
\langle|\Delta\kappa|\rangle_{\rm near}.
\]

Avec \(R^2=0.967\) sur le tore plat et \(R^2=0.982\) sur la sphère, la réponse
modulaire prédit fortement la réponse de courbure proche.

Ainsi, Paper~18 fournit la brique source du programme BuP :
\[
\delta W_{\rm loc}
\longrightarrow
T_{\mu\nu}^{\rm ent}.
\]

Il ne s'agit pas encore d'une dérivation tensorielle complète, mais d'une
validation numérique contrôlée montrant que la variation modulaire se comporte
comme une source de courbure.

\section*{Disponibilité du code et des données}

Les scripts et résultats de Paper~18 sont organisés dans le dossier :
\[
\texttt{papers/paper18\_entanglement\_stress\_tensor/}.
\]

Les scripts principaux sont :
\[
\texttt{scripts/paper18\_graph\_modular\_first\_law\_v1.py},
\]
\[
\texttt{scripts/paper18\_modular\_source\_curvature\_v2.py},
\]
\[
\texttt{scripts/paper18\_build\_entanglement\_stress\_summary\_v1.py}.
\]

Les résultats principaux sont stockés dans :
\[
\texttt{results/graph\_modular\_first\_law\_v1/},
\]
\[
\texttt{results/modular\_source\_curvature\_v2/},
\]
\[
\texttt{results/paper18\_entanglement\_stress\_summary\_v1/}.
\]

\appendix

\section{Prochaines étapes}

Les prochaines étapes sont :

\begin{enumerate}
    \item reconstruire une vraie matrice densité réduite \(\rho_A\) depuis un état quantique ;
    \item comparer le proxy graphe à \(K_A=-\log\rho_A\) exact ;
    \item tester des sources anisotropes pour reconstruire plusieurs composantes de \(T_{\mu\nu}\) ;
    \item relier \(\delta\langle K_A\rangle\) à un courant d'échange d'intrication ;
    \item formaliser la limite continue via la première loi de l'intrication ;
    \item intégrer ce résultat dans Paper~19, l'assemblage de l'équation d'Einstein effective.
\end{enumerate}

\end{document}
